\documentclass{article}
\usepackage{graphicx}

\usepackage[utf8]{inputenc}     % pdfLaTeX
\usepackage[T1]{fontenc}
\usepackage[magyar]{babel}

\title{Az Idegrendszer}
\author{Lengyel Zoltán Péter, Rovenszky Robin}
\date{Legutóbb frissítve: 2026. Április 22.}

\begin{document}

\section{Az idegrendszer működése}

\section{Az idegrendszer felépítése}
\begin{itemize}
    \item Központi idegrendszer \begin{itemize}
        \item Agyvelő
        \item Gerincvelő
    \end{itemize}
\end{itemize}

\subsection{Központi idegrendszer}
\subsubsection{Agyvelő}
\begin{itemize}
    \item agytörzs \begin{enumerate}
        \item nyúltagy (gerincvelőhöz kapcsolódik)
        \item híd (kiszélesedés)
        \item középagy (híd fölött)
    \end{enumerate}
    \item kisagy
    \item köztiagy
    \item nagyagy
\end{itemize}

\subsubsection{Gerincvelő}
\begin{itemize}
    \item Szürkeállomány \begin{itemize}
        \item Sejttestek tömege (mozgató neuronok, asszociációs neuronok)
        \item Hátsó szarv (hát fele) \begin{itemize}
            \item ide lépnek be az érző neuronok, belépés előtt érződúc (érzőneuron sejttestje, pseudo-unipoláris sejt, ál-kétnyúlványú idegsejt)
        \end{itemize}
        \item Mellső szarv (has fele)
    \end{itemize}
    \item Fehérállomány \begin{itemize}
        \item Agyba felszálló pályák és agyból jövő pályák tömege
    \end{itemize}
    \item Központi csatorna \begin{itemize}
        \item liquor [ejtsd: likuor] táplálja az idegsejteket, zárt, steril rendszer
    \end{itemize}
\end{itemize}

% insert "az ember idegrendszere" here from ppt
\begin{tabular}{|c|c|c|}
    \hline
     Központi idegrendszer & szürkeállomány & fehérállomány \\
     \hline
     & kéreg és magvak & pályák \\
     & sejttest+dendritek & axon \\
     \hline
     Környéki idegrendszer & dúc & ideg \\
     \hline
\end{tabular}



% insert "Az idegrendszer felosztása" here
\end{document}